% Full Etemadi and Kronecker/random-series routes, with their common estimates.
\section{Strong laws of large numbers: proofs}
\label{sec:annex-p15-strong-laws-of-large-numbers}

Use Chapter~\ref{sec:p15-strong-laws-of-large-numbers} for all statements
and definitions. Durrett's Theorem 2.4.1 is proved here through its full
truncation, geometric-subsequence and interpolation argument. The
distinct proof of Theorem 2.5.10 follows Kronecker's lemma and the
independent-series convergence theorem
\cite[pp.~76--78, 85--87]{DurrettPTEAuthor2019}.

\begin{annexitem}{annex:p15.cesaro-average}{thm:p15-cesaro-average}{The finite head and small tail of an average}
\begin{proof}
Given \(\varepsilon>0\), choose \(K\geq1\) such that
\(|u_k-u|\leq\varepsilon\) for every \(k\geq K\). For
\(n\geq K\),
\[
 \left|\frac1n\sum_{k=1}^nu_k-u\right|
 \leq\frac1n\sum_{k=1}^{K-1}|u_k-u|
                          +\frac{n-K+1}{n}\varepsilon.
\]
The first numerator is a fixed finite number and the second coefficient
is at most one. The limit superior of the left side is therefore at
most \(\varepsilon\). As \(\varepsilon\) is arbitrary and
the left side is nonnegative, it tends to zero.
\end{proof}
\end{annexitem}

\begin{annexitem}{annex:p15.strong-law-tail-bounds}{thm:p15-strong-law-tail-bounds}{Tail integration and both ranges of the numerical bound}
\begin{proof}
The first identity is Theorem~\ref{thm:p15-layer-cake} on the
finite probability space. For the second, at every outcome,
\[
 W(\omega)^2=\int_0^\infty
                       2t\mathbf1_{\{t<W(\omega)\}}\,dt.
\]
This is the elementary integral of \(2t\) over \([0,W(\omega))\).
The subgraph indicator is product measurable by the same layer-cake
theorem, and multiplication by the nonnegative Borel function \(2t\)
preserves measurability. The probability factor is finite and Lebesgue
measure on \([0,\infty)\) is sigma-finite. Tonelli's theorem therefore
interchanges the two nonnegative integrals, proving
\(\mathbb EW^2=\int_0^\infty2t\mathbb P(W>t)\,dt\), even
when this value is infinite. No moment finiteness or signed integral
interchange is assumed in this step.

For \(t\in[k-1,k)\), the event \(\{W>k\}\) is contained
in \(\{W>t\}\). Integrating over these disjoint unit intervals
and summing gives
\[
 \sum_{k\geq1}\mathbb P(W>k)
  \leq\sum_{k\geq1}\int_{k-1}^k\mathbb P(W>t)\,dt
  =\mathbb EW.
\]
The equality uses countable additivity for nonnegative integrals.

For each integer \(m\geq2\), monotonicity of \(x^{-2}\) on
\((0,\infty)\) yields
\[
 \sum_{k=m}^\infty k^{-2}
     \leq\sum_{k=m}^\infty\int_{k-1}^k x^{-2}\,dx
     =\int_{m-1}^\infty x^{-2}\,dx
     =\frac1{m-1}.
\]
If \(y\geq1\), put \(r=\lfloor y\rfloor\geq1\). Then
\(\sum_{k>y}k^{-2}=\sum_{k\geq r+1}k^{-2}\leq1/r\).
Since \(y<r+1\leq2r\), multiplication by \(2y\) gives the
bound at most four. If \(0\leq y<1\), the sum includes all
positive integers, and
\[
 2y\sum_{k>y}k^{-2}
   \leq2\left(1+\sum_{k=2}^\infty k^{-2}\right)\leq4.
\]
At \(y=0\) the left side is zero; the series multiplying it is
finite by the displayed estimate. This proves every stated range.
\end{proof}
\end{annexitem}

\begin{annexitem}{annex:p15.strong-law-truncation}{thm:p15-strong-law-truncation}{Finite modifications, summable scaled variances and common-law means}
\begin{proof}
All \(Y_k\) are measurable and bounded in absolute value by \(k\),
so each has finite first and second moments. The event \(X_k\ne Y_k\)
is exactly \(\{|X_k|>k\}\). Identical laws and the preceding
tail-probability bound give
\[
 \sum_{k\geq1}\mathbb P(X_k\ne Y_k)
       =\sum_{k\geq1}\mathbb P(|X_1|>k)
       \leq\mathbb E|X_1|<\infty.
\]
The first Borel--Cantelli lemma implies that only finitely many changes
occur almost surely; independence is unnecessary. The measurable
extended nonnegative variable
\[
 R=\sum_{k\geq1}|X_k|\mathbf1_{\{|X_k|>k\}}
\]
is finite on this probability-one event, since only finitely many
finite summands are nonzero there. At every outcome and every \(n\),
\(|S_n-T_n|\leq R\). In particular, on the event where \(R<\infty\),
division by \(n\) proves \((S_n-T_n)/n\to0\). No integrability
claim about the random bound \(R\) is needed.

For the variance estimate, \(\operatorname{Var}(Y_k)\leq\mathbb EY_k^2\).
If \(y\geq0\), then
\[
 \mathbb P(|Y_k|>y)
   \leq\mathbf1_{\{y<k\}}\mathbb P(|X_1|>y).
\]
Indeed the left side is zero for \(y\geq k\), and for \(y<k\)
its event implies \(|X_k|>y\), which has the same probability
as \(|X_1|>y\). Applying the second tail-integral formula and
interchanging the nonnegative sum and integral gives
\[
 \begin{aligned}
 \sum_{k\geq1}\frac{\mathbb EY_k^2}{k^2}
 &\leq\sum_{k\geq1}k^{-2}
             \int_0^\infty2y\mathbf1_{\{y<k\}}
                                      \mathbb P(|X_1|>y)\,dy\\
 &=\int_0^\infty
       \left(2y\sum_{k>y}k^{-2}\right)\mathbb P(|X_1|>y)\,dy\\
 &\leq4\int_0^\infty\mathbb P(|X_1|>y)\,dy
       =4\mathbb E|X_1|.
 \end{aligned}
\]
The interchange is Tonelli on positive integer counting measure times
Borel Lebesgue measure, both sigma-finite. The integrand is nonnegative
and product measurable; equivalently monotone convergence applies to
the finite sums. The numerical bound in
Lemma~\ref{thm:p15-strong-law-tail-bounds} proves the last inequality.

For the means, identical distributions give the equality
\[
 \mathbb EY_k
     =\mathbb E\!\left[X_1\mathbf1_{\{|X_1|\leq k\}}\right]
\]
by integration over the common law
(Theorem~\ref{thm:p15-expectation-law-integral}). The functions of
\(X_1\) on the right tend pointwise to \(X_1\), because it is
real valued, and are dominated in absolute value by the integrable
\(|X_1|\). Dominated convergence proves \(\mathbb EY_k\to\mu\).
This is a common-law calculation; it does not assert that the actual
sequence \(Y_k(\omega)\) converges to \(X_1(\omega)\).
Finite linearity and Lemma~\ref{thm:p15-cesaro-average} then give
\(\mathbb ET_n/n=n^{-1}\sum_{k=1}^n\mathbb EY_k\to\mu\).

Finally, for each \(k\), deletion truncation is a Borel function
of its single input. Any two independent inputs give independent
outputs by the generated-independence results. This proves preservation
of pairwise independence. The same argument applied to each finite
subfamily proves preservation of mutual independence when present.
\end{proof}
\end{annexitem}

\begin{annexitem}{annex:p15.geometric-subsequence}{thm:p15-geometric-subsequence}{The floor estimate, sum interchange and countable thresholds}
\begin{proof}
For every real \(x\geq1\), \(\lfloor x\rfloor\geq x/2\):
if \(1\leq x<2\), use \(\lfloor x\rfloor=1\); if \(x\geq2\),
use \(\lfloor x\rfloor>x-1\geq x/2\), with the integer endpoint
cases satisfying the same inequality. Also
\(0\leq\alpha^j-k(j)<1\). Hence
\(k(j)/\alpha^j\to1\), proving divergence and the consecutive
ratio limit. The indices are nondecreasing; a few initial values can
repeat, which does not affect any of these conclusions.

For a fixed positive integer \(m\), let \(j_0\geq1\) be the
least integer for which \(\alpha^{j_0}\geq m\). If
\(k(j)\geq m\), then \(\alpha^j\geq m\). Therefore
\[
 \begin{aligned}
 \sum_{j:k(j)\geq m}k(j)^{-2}
 &\leq4\sum_{j:\alpha^j\geq m}\alpha^{-2j}\\
 &=\frac{4\alpha^{-2j_0}}{1-\alpha^{-2}}
 \leq\frac{4m^{-2}}{1-\alpha^{-2}}.
 \end{aligned}
\]
The denominator is positive because \(\alpha>1\).

For the random-variable assertion, pairwise independence and finite
second moments imply
\(\operatorname{Var}(T_\ell)=\sum_{m=1}^\ell\operatorname{Var}(Y_m)\)
by Theorem~\ref{thm:p15-independent-moments}. Chebyshev's inequality
thus gives
\[
 \begin{aligned}
 &\sum_{j\geq1}\mathbb P\!\left(
     |T_{k(j)}-\mathbb ET_{k(j)}|>\varepsilon k(j)\right)\\
 &\quad\leq\varepsilon^{-2}\sum_{j\geq1}k(j)^{-2}
                        \sum_{m=1}^{k(j)}\operatorname{Var}(Y_m)\\
 &\quad=\varepsilon^{-2}\sum_{m\geq1}\operatorname{Var}(Y_m)
                               \sum_{j:k(j)\geq m}k(j)^{-2}\\
 &\quad\leq\frac{4\varepsilon^{-2}}{1-\alpha^{-2}}
                   \sum_{m\geq1}\frac{\operatorname{Var}(Y_m)}{m^2}.
 \end{aligned}
\]
All terms are nonnegative, so the double sum can be interchanged by
Tonelli for the two counting measures. The right side is finite by
hypothesis. Apply the first Borel--Cantelli lemma at each threshold
\(\varepsilon=1/r\), \(r\geq1\), and intersect the resulting
countably many probability-one events. At an outcome in this
intersection, the normalized centered deviations are eventually at
most every \(1/r\) in absolute value, hence tend to zero. No
independence between the deviation events is needed.
\end{proof}
\end{annexitem}

\begin{annexitem}{annex:p15.etemadi-strong-law}{thm:p15-etemadi-strong-law}{The full pairwise proof and interpolation between geometric indices}
\begin{proof}
We first prove the result for nonnegative \(X_n\), and then reduce
the general real case to two such sequences. In the nonnegative case
the moving truncations \(Y_n\) are nonnegative, so their partial
sums \(T_n\) are nondecreasing. All have finite second moments,
and they remain pairwise independent. The truncation theorem gives
\(\sum_n\operatorname{Var}(Y_n)/n^2<\infty\) and
\(\mathbb ET_n/n\to\mu=\mathbb EX_1\geq0\).

Fix \(\alpha>1\) and use \(k(j)=\lfloor\alpha^j\rfloor\).
Lemma~\ref{thm:p15-geometric-subsequence} proves almost surely
\[
 \frac{T_{k(j)}-\mathbb ET_{k(j)}}{k(j)}\longrightarrow0.
\]
Adding the convergent deterministic expected averages gives
\(T_{k(j)}/k(j)\to\mu\) on an event \(E_\alpha\) of
probability one.

For each integer \(m\geq k(1)\), choose the largest integer
\(j\) with \(k(j)\leq m\). This integer exists since the
indices are unbounded; using the largest handles any repeated initial
values and ensures \(k(j)\leq m<k(j+1)\). The chosen \(j\)
tends to infinity as \(m\to\infty\). Nonnegativity and
monotonicity imply
\[
 \frac{T_{k(j)}}{k(j+1)}
       \leq\frac{T_m}{m}
       \leq\frac{T_{k(j+1)}}{k(j)}.
\]
On \(E_\alpha\), the left side is
\((T_{k(j)}/k(j))(k(j)/k(j+1))\to\mu/\alpha\),
and the right side is
\((T_{k(j+1)}/k(j+1))(k(j+1)/k(j))\to\alpha\mu\).
Consequently
\[
 \mu/\alpha\leq\liminf_{m\to\infty}T_m/m
       \leq\limsup_{m\to\infty}T_m/m\leq\alpha\mu
                                      \quad\text{on }E_\alpha.
\]

To let the ratio tend to one on a common event, take only the countable
sequence \(\alpha_\ell=1+1/\ell\), \(\ell\geq1\).
The intersection \(E=\bigcap_{\ell\geq1}E_{\alpha_\ell}\)
has probability one. At any outcome in \(E\), the last bounds
hold for every \(\ell\). Letting \(\ell\to\infty\) gives
\(\liminf T_m/m\geq\mu\) and \(\limsup T_m/m\leq\mu\).
They coincide, including when \(\mu=0\). Thus \(T_m/m\to\mu\)
almost surely. The finite-change estimate
\((S_m-T_m)/m\to0\) proves the same limit for \(S_m/m\).
This establishes the nonnegative case using only pairwise independence.

For real integrable \(X_n\), define
\(X_n^+=\max(X_n,0)\) and \(X_n^-=\max(-X_n,0)\).
Each sequence separately is pairwise independent, since each output
is a measurable function of one input, and identically distributed,
since the same function is applied at every index. Both are integrable
because \(0\leq X_n^\pm\leq|X_n|\). Apply the nonnegative case
just proved to each sequence. On the intersection of its two
probability-one events,
\[
 \frac1n\sum_{k=1}^nX_k
   =\frac1n\sum_{k=1}^nX_k^+
     -\frac1n\sum_{k=1}^nX_k^-
   \longrightarrow\mathbb EX_1^+-\mathbb EX_1^-=\mathbb EX_1.
\]
Both limiting expectations are finite, so their subtraction is
defined. No independence between \(X_n^+\) and \(X_n^-\) is
asserted or needed. This completes the full signed theorem.
\end{proof}
\end{annexitem}

\begin{annexitem}{annex:p15.kronecker-lemma}{thm:p15-kronecker-lemma}{Summation by parts and concentration of the deterministic weights}
\begin{proof}
First take the positive nondecreasing weights. Let \(b_n\) be the
partial sums of \(\sum_kx_k/a_k\), and let \(b=\lim_nb_n\)
be finite. Since \(x_m=a_m(b_m-b_{m-1})\), finite reindexing gives
\[
 \begin{aligned}
 \sum_{m=1}^nx_m
  &=\sum_{m=1}^na_mb_m-\sum_{m=1}^na_mb_{m-1}\\
  &=a_nb_n+\sum_{m=2}^na_{m-1}b_{m-1}
                                  -\sum_{m=1}^na_mb_{m-1}\\
  &=a_nb_n-\sum_{m=1}^n(a_m-a_{m-1})b_{m-1}.
 \end{aligned}
\]
The last equality uses \(a_0b_0=0\); the sums with no terms
are zero when \(n=1\). Division by the positive \(a_n\) gives
the stated identity.

Put \(q_{n,m}=(a_m-a_{m-1})/a_n\), for \(1\leq m\leq n\).
These numbers are nonnegative, including at \(m=1\), and
\(\sum_{m=1}^nq_{n,m}=(a_n-a_0)/a_n=1\).
For fixed \(M\leq n\), their total mass on \(m\leq M\)
is \(a_M/a_n\), which tends to zero as \(n\to\infty\).
To prove the required weighted limit explicitly, let \(\varepsilon>0\)
and choose \(M\geq1\) such that \(|b_j-b|\leq\varepsilon\)
for all \(j\geq M\). The finite number
\(H_M=\max_{0\leq j<M}|b_j-b|\) bounds the earlier errors.
For \(n\geq M\),
\[
 \begin{aligned}
 \left|\sum_{m=1}^nq_{n,m}b_{m-1}-b\right|
 &\leq\sum_{m=1}^nq_{n,m}|b_{m-1}-b|\\
 &\leq\frac{a_M}{a_n}H_M
             +\frac{a_n-a_M}{a_n}\varepsilon.
 \end{aligned}
\]
Thus the limit superior of this nonnegative error is at most
\(\varepsilon\). Letting \(\varepsilon\downarrow0\) proves
\(\sum_mq_{n,m}b_{m-1}\to b\). This estimate also covers
\(H_M=0\) and the identically zero partial-sum sequence; no division
by an auxiliary error bound is needed. Since \(b_n\to b\),
the summation identity proves \(a_n^{-1}\sum_{k=1}^nx_k\to0\).

For the extended weight domain, choose \(N\) after which all weights
are positive and nondecreasing. Removing the first \(N-1\) terms
from the convergent series leaves a convergent tail. Apply the positive
case to \(\widetilde a_j=a_{N+j-1}\) and
\(\widetilde x_j=x_{N+j-1}\). It gives
\(a_n^{-1}\sum_{k=N}^nx_k\to0\). The omitted finite sum
\(\sum_{k=1}^{N-1}x_k\), divided by \(a_n\to\infty\), also
tends to zero. Adding it proves the full conclusion. The requirement
\(a_n\ne0\) ensures that every term of the original weighted series
is defined.
\end{proof}
\end{annexitem}

\begin{annexitem}{annex:p15.weighted-independent-strong-law}{thm:p15-weighted-independent-strong-law}{Centering, scaling and the pointwise Kronecker argument}
\begin{proof}
Let \(Z_n=(X_n-\mu_n)/a_n\). These real variables are mutually
independent measurable functions of their separate inputs. They are
centered and have finite variances
\(\operatorname{Var}(Z_n)=\operatorname{Var}(X_n)/a_n^2\).
Their variance series is finite by hypothesis, so
Theorem~\ref{thm:p15-variance-series-convergence} gives almost-sure
finite convergence of \(\sum_nZ_n\).
At every outcome of that probability-one convergence event, apply
Kronecker's deterministic lemma with
\(x_n=X_n(\omega)-\mu_n\). Its weights are the given \(a_n\),
and its convergent weighted series is exactly \(\sum_nZ_n(\omega)\).
It follows on that event that
\(a_n^{-1}\sum_{k=1}^n(X_k-\mu_k)\to0\).
For \(a_n=n\), add the assumed convergent numerical averages of
the means to obtain \(\overline X_n\to\mu\).
\end{proof}
\end{annexitem}

\begin{annexitem}{annex:p15.iid-series-strong-law}{thm:p15-iid-series-strong-law}{The full independent-series proof of the integrable law}
\begin{proof}
We use the common-law truncation estimates but do not invoke the
pairwise strong-law theorem. Put
\(Y_k=X_k\mathbf1_{\{|X_k|\leq k\}}\),
\(T_n=\sum_{k=1}^nY_k\), and \(Z_k=Y_k-\mathbb EY_k\).
Since the original sequence is mutually independent, so are the
\(Y_k\), the \(Z_k\), and the scaled variables \(Z_k/k\).
Every \(Z_k/k\) is centered with finite second moment. The
previously proved truncation bounds give
\[
 \sum_{k\geq1}\operatorname{Var}(Z_k/k)
  =\sum_{k\geq1}\frac{\operatorname{Var}(Y_k)}{k^2}
  \leq\sum_{k\geq1}\frac{\mathbb EY_k^2}{k^2}
  \leq4\mathbb E|X_1|<\infty.
\]
The independent variance-series convergence theorem gives almost-sure
finite convergence of \(\sum_k Z_k/k\). Apply Kronecker's lemma
pointwise on its probability-one convergence event, with weights
\(a_k=k\). Then
\[
 \frac1n\sum_{k=1}^n(Y_k-\mathbb EY_k)
   =\frac{T_n}{n}-\frac1n\sum_{k=1}^n\mathbb EY_k
                                  \longrightarrow0.
\]
The common-law dominated-convergence calculation in
Theorem~\ref{thm:p15-strong-law-truncation} gives
\(\mathbb EY_k\to\mu\), and the numerical averaging lemma
therefore gives \(n^{-1}\sum_{k=1}^n\mathbb EY_k\to\mu\).
Adding the last two limits proves \(T_n/n\to\mu\) almost surely.
The same truncation theorem gives \((S_n-T_n)/n\to0\) almost
surely. On the intersection of these probability-one events we obtain
\(S_n/n\to\mu\), as required.

Only the truncated variables needed finite second moments in this
argument; the original \(X_k\) were assumed merely integrable.
Mutual independence was used at the invocation of the independent-series
theorem, whose proof controls maxima of sums through past/future
independence. The pairwise proof above reaches its conclusion through
different estimates and does not make that invocation.
\end{proof}
\end{annexitem}

\begin{annexitem}{annex:p15.strong-law-applications}{thm:p15-strong-law-applications}{Measurable observations and the computed integrals}
\begin{proof}
For a fixed Borel \(g\), its outputs \(g(X_n)\) have the same
law at every index and remain pairwise independent. Integration over
the common input law gives
\[
 \mathbb E|g(X_n)|=\int|g|\,d\nu<\infty,
 \qquad \mathbb Eg(X_n)=\int g\,d\nu.
\]
The pairwise strong law applied to this output sequence proves the
claimed empirical-average limit. For a specified countable family
\(g_1,g_2,\ldots\), intersect the associated countably many
probability-one events. This gives simultaneous convergence on their
intersection, without making an uncountable-intersection assertion.
Taking \(g=\mathbf1_B\) for a Borel set \(B\) gives the frequency
limit, since the indicator is bounded and its integral is \(\nu(B)\).

For the uniform law, its restriction of Lebesgue measure to \([0,1]\)
has total mass one. The measurable function \(u^2\) is bounded
there, so the integrability condition holds. The elementary integral
\(\int_0^1u^2\,du=[u^3/3]_0^1=1/3\) gives the first value.
For the second example the proposed density is nonnegative and Borel,
and its integral is
\[
 \int_1^\infty\frac32x^{-5/2}\,dx
      =\lim_{b\to\infty}(1-b^{-3/2})=1.
\]
It therefore defines a probability law. Its first absolute moment
equals its first moment, since its support is positive, and
\[
 \int_1^\infty x\frac32x^{-5/2}\,dx
      =\lim_{b\to\infty}3(1-b^{-1/2})=3.
\]
Its second moment is
\[
 \int_1^\infty x^2\frac32x^{-5/2}\,dx
      =\lim_{b\to\infty}3(\sqrt b-1)=\infty.
\]
All three computations use nonnegative integrands and monotone
limits of bounded-interval integrals, so they do not assume integrability
before establishing it. Independent sequences of either specified
probability law exist by
Theorem~\ref{thm:p15-independent-sequence-existence}. Applying
either integrable strong law gives the asserted averages.
\end{proof}
\end{annexitem}

\begin{annexitem}{annex:p15.pairwise-block-example}{thm:p15-pairwise-block-example}{Independent blocks with a deterministic triple constraint}
\begin{proof}
The independent signs \((U_j,V_j)_{j\geq1}\) can be constructed
by the countable independent-sequence theorem, using alternating
coordinates for the two signs in each pair. Within a fixed pair
\((U_j,V_j)\), each of the four points \((u,v)\in\{-1,1\}^2\)
has probability \(1/4\). The product \(U_jV_j\) is consequently
a fair sign. Moreover, the maps
\((u,v)\mapsto(u,uv)\) and \((u,v)\mapsto(v,uv)\) are
bijections of this four-point set. Thus the pair
\((U_j,U_jV_j)\), and likewise \((V_j,U_jV_j)\), has uniform
joint law on four points and is independent. The remaining pair
\((U_j,V_j)\) is independent by construction.

Outputs from different triples are functions of disjoint independent
input pairs, so every such output pair is independent as well.
This proves pairwise independence of the whole \(X_n\) sequence,
and every marginal is the fair-sign law. But
\(X_{3j-2}X_{3j-1}X_{3j}=U_jV_j(U_jV_j)=1\) at every outcome.
All three outputs are positive exactly when \(U_j=V_j=1\),
which has probability \(1/4\). If the triple were mutually
independent, that probability would instead be \((1/2)^3=1/8\).
Thus mutual independence fails. The variables are bounded and have
mean zero, so the pairwise strong law applies and proves the stated
zero limit of their sample averages.
\end{proof}
\end{annexitem}

\begin{annexitem}{annex:p15.strong-law-counterexamples}{thm:p15-strong-law-counterexamples}{Explicit failures of the distribution and weight hypotheses}
\begin{proof}
For the shared-sign example, every variable has the fair-sign law,
so its mean is zero and its absolute first moment is one. The events
\(\{X_1=1\}\) and \(\{X_2=1\}\) coincide and have probability
\(1/2\), whereas the product of their probabilities is \(1/4\).
They are not independent. Direct summation gives
\(S_n=n\varepsilon\) and \(S_n/n=\varepsilon\), which never
equals the asserted mean-zero limit.

For the independent but nonidentical example, the specified masses
are nonnegative and sum to one at every index, including \(n=1\).
Symmetry gives \(\mathbb EX_n=0\), and
\(\mathbb E|X_n|=n(1/n)=1\). The independent events
\(A_n=\{|X_n|=n\}\) have probabilities \(1/n\), whose sum
is infinite. For example, each dyadic block
\(2^r\leq n<2^{r+1}\) contributes at least \(1/2\).
The second Borel--Cantelli lemma gives \(A_n\) infinitely often
almost surely. Hence \(X_n/n\) does not tend to zero on that
probability-one event.
On the other hand, if at an outcome \(S_n/n\to L\in\mathbb R\),
then for \(n\geq2\),
\[
 \frac{X_n}{n}
   =\frac{S_n}{n}-\frac{n-1}{n}\frac{S_{n-1}}{n-1}
                           \longrightarrow L-L=0.
\]
Thus the event of a finite sample-average limit is contained in the
null complement of the infinitely-often event. This proves almost-sure
failure even though all the means, and their averages, are zero.

For the positive dyadic common law,
\(\sum_{j=1}^\infty2^{-j}=1\), so the masses define a probability
law on finite positive real values. Its first moment is
\(\sum_{j=1}^\infty2^j2^{-j}=\sum_{j=1}^\infty1=\infty\).
For each integer \(r\geq1\), use the bounded clipped observation
\(X_n\wedge2^r=\min(X_n,2^r)\). These variables are i.i.d.
and integrable; their mean is
\[
 \mathbb E(X_1\wedge2^r)
  =\sum_{j=1}^r2^j2^{-j}
       +2^r\sum_{j=r+1}^\infty2^{-j}=r+1.
\]
The integrable strong law just proved implies that their averages
tend to \(r+1\) on an event of probability one. Intersect these
events over the positive integers \(r\). On the intersection,
\(X_n\geq X_n\wedge2^r\) for every \(n,r\), so
\[
 \liminf_{n\to\infty}S_n/n\geq r+1
                                  \quad\text{for every }r\geq1.
\]
The limit inferior is therefore infinite, which is exactly
\(S_n/n\to\infty\). This establishes the displayed example
using bounded observations of that law; it does not assume an
infinite-mean strong law as a prerequisite.

For the nonmonotone Kronecker weights, both subsequences of weights
tend to infinity, so the full positive sequence does too. Its successive
odd-to-even drops occur for every \(j\geq2\), since \(j^3>j\),
so it is not eventually nondecreasing. The weighted summands are
\[
 x_{2j-1}/a_{2j-1}=1/j^2,\qquad
 x_{2j}/a_{2j}=-1/j^2.
\]
Their absolute series has sum \(2\sum_jj^{-2}<\infty\), using
the reciprocal-square estimate already proved. Writing
\(H_j=\sum_{m=1}^j1/m\), direct summation of the unweighted terms
gives
\[
 \frac1{a_{2j}}\sum_{k=1}^{2j}x_k
   =\frac1j\left(\frac{j(j+1)}2-H_j\right)
   =\frac{j+1}{2}-\frac{H_j}{j}.
\]
Since \(H_j\leq j\), this is at least \((j-1)/2\), which
tends to infinity. Finally, with constant weights \(a_n=1\)
and the single nonzero term \(x_1=1\), the weighted series
converges to one but every normalized unweighted partial sum is one.
These computations prove the claimed failures of the two weight
hypotheses.
\end{proof}
\end{annexitem}
